% ============================================================
%  GameNight - Advanced DBMS Semester Project Report
% ============================================================
\documentclass[11pt,a4paper]{article}

\usepackage[utf8]{inputenc}
\usepackage[T1]{fontenc}
\usepackage[margin=1in]{geometry}
\usepackage{graphicx}
\usepackage{booktabs}
\usepackage{longtable}
\usepackage{xcolor}
\usepackage{listings}
\usepackage{enumitem}
\usepackage{float}
\usepackage{titlesec}
\usepackage[hidelinks]{hyperref}

% ---- Colors ----
\definecolor{accent}{HTML}{7C3AED}
\definecolor{codebg}{HTML}{F5F5FA}
\definecolor{codekw}{HTML}{6D28D9}
\definecolor{codecm}{HTML}{6B7280}
\definecolor{codestr}{HTML}{0E7490}

% ---- Listings (SQL / JS / shell) ----
\lstdefinestyle{base}{
  backgroundcolor=\color{codebg},
  basicstyle=\ttfamily\small,
  keywordstyle=\color{codekw}\bfseries,
  commentstyle=\color{codecm}\itshape,
  stringstyle=\color{codestr},
  breaklines=true,
  showstringspaces=false,
  frame=single,
  rulecolor=\color{gray!30},
  framesep=6pt,
  xleftmargin=8pt,
  numbers=none,
  captionpos=b,
}
\lstset{style=base}

% ---- Section styling ----
\titleformat{\section}{\Large\bfseries\color{accent}}{\thesection}{1em}{}
\titleformat{\subsection}{\large\bfseries}{\thesubsection}{1em}{}

% ---- Screenshot helper: shows image if present (in ../screenshots or
%      ./screenshots), else a placeholder box so the document still compiles ----
\newcommand{\screenshot}[3]{%
  \begin{figure}[H]
    \centering
    \IfFileExists{../screenshots/#1}{%
      \fbox{\includegraphics[width=0.85\textwidth]{../screenshots/#1}}%
    }{%
      \IfFileExists{screenshots/#1}{%
        \fbox{\includegraphics[width=0.85\textwidth]{screenshots/#1}}%
      }{%
        \fbox{\parbox[c][4cm][c]{0.8\textwidth}{\centering\color{gray}\itshape
          [ screenshot placeholder ]\\\texttt{screenshots/#1}}}%
      }%
    }
    \caption{#2}\label{#3}
  \end{figure}%
}

\begin{document}

% ===================== TITLE =====================
\begin{titlepage}
\centering
\vspace*{2cm}
{\Huge\bfseries\color{accent} GameNight\par}
\vspace{0.4cm}
{\Large A Group Video Game Night Planner\par}
\vspace{1.5cm}
{\large\bfseries Advanced Database Management Systems\\Individual Semester Project\par}
\vspace{2cm}
\begin{tabular}{rl}
\textbf{Student:} & Muhammad Izzat Kamal \\[4pt]
\textbf{Roll No:} & BSAI24078 \\[4pt]
\textbf{Course:} & Advanced Database Management Systems \\
\end{tabular}
\vfill
\begin{tabular}{rl}
\textbf{Live App:} & \url{https://gamenight-adbms.vercel.app} \\[4pt]
\textbf{Repository:} & \url{https://github.com/izzatkamall/gamenight-adbms} \\
\end{tabular}
\vspace{1cm}
\end{titlepage}

\tableofcontents
\newpage

% ===================== 1. DESCRIPTION =====================
\section{Project Title and Description}
\textbf{GameNight} is a database-backed web application that helps friend groups
decide which video game to play together. It tackles the everyday problem of
conflicting game libraries, differing genre tastes, and group indecision by:

\begin{itemize}[leftmargin=1.4em]
  \item Computing the set of games that \emph{every} member of a room owns in common.
  \item Building per-user taste profiles from session ratings, stored as JSONB.
  \item Ranking a shortlist of candidate games by aggregate group preference.
  \item Running a \emph{live, real-time vote} among the group to pick the winner.
  \item Logging every game-night session so future recommendations improve.
\end{itemize}

The project is intentionally designed to showcase a broad range of advanced
database concepts within a single, coherent application --- JSONB document
modelling inside a relational database, set-theoretic intersection queries,
stored procedures, triggers, row-level security, an in-memory store (Redis)
for high-frequency writes, and real-time subscriptions.

% ===================== 2. OBJECTIVES & SCOPE =====================
\section{Objectives and Scope}
\subsection{Objectives}
\begin{itemize}[leftmargin=1.4em]
  \item Design a normalised (3NF) relational schema for users, games, libraries,
        rooms, voting, and session history.
  \item Demonstrate the integration of a NoSQL-style JSONB document model
        \emph{within} PostgreSQL for flexible preference profiles.
  \item Implement core business logic as PostgreSQL stored functions rather than
        in application code.
  \item Use Redis as a justified in-memory store for the write-heavy, latency-sensitive
        live-voting feature.
  \item Enforce data isolation between users and rooms using Row-Level Security (RLS).
  \item Deploy the application so it is publicly accessible.
\end{itemize}

\subsection{Scope}
\textbf{In scope:} user registration/authentication, personal game libraries,
room creation and joining, common-library computation, preference profiles,
ranked shortlists, real-time voting, session logging and ratings, session
history, and a set of bonus DBMS features (materialized views, window functions,
full-text search, audit logging via triggers, real-time chat, and wishlists).

\textbf{Out of scope:} integration with external game-store APIs (e.g.\ Steam),
payment processing, mobile-native apps, and matchmaking with strangers. The game
catalogue is a curated seed set rather than a live external feed.

% ===================== 3. TECH STACK & ARCHITECTURE =====================
\section{Tech Stack and Architecture}
\subsection{Technology Stack}
\begin{longtable}{p{3.2cm} p{10.5cm}}
\toprule
\textbf{Layer} & \textbf{Technology} \\
\midrule
Frontend        & React + Vite + Tailwind CSS (dark/glassmorphism UI) \\
Primary Database & PostgreSQL via Supabase --- JSONB, stored functions, triggers, RLS, GIN indexes \\
Authentication  & Supabase Auth (email/password) \\
Real-time       & Supabase Realtime (\texttt{postgres\_changes}) for room state and chat \\
In-memory Store & Redis --- Sorted Sets and Sets for live voting \\
Middleware      & Node.js + Express + WebSocket (\texttt{ws}) + \texttt{ioredis} \\
\bottomrule
\end{longtable}

\subsection{Architecture Overview}
The frontend talks \emph{directly} to Supabase for all standard CRUD, auth, and
real-time subscriptions. A separate lightweight Node.js middleware exists for one
purpose only: bridging the browser to Redis over WebSockets for the live vote,
where atomic, sub-millisecond increments are required.

\begin{lstlisting}[basicstyle=\ttfamily\scriptsize]
+-------------+   Supabase JS client (auth, queries,   +----------------------+
|  React App  |<--  real-time subscriptions) -------->  |  Supabase            |
|  (Vite)     |                                         |  PostgreSQL +        |
|             |                                         |  Realtime + Auth     |
|             |   WebSocket                             +----------------------+
|             |<----------------------------->  +----------------------+
+-------------+                                 |  Node.js middleware   |
                                                |  (Express + ws)       |
                                                +----------+-----------+
                                                           | Redis commands
                                                +----------v-----------+
                                                |  Redis (Sorted Sets) |
                                                +----------------------+
\end{lstlisting}

\textbf{Live-voting data flow:} (1) React opens a WebSocket to the middleware;
(2) the middleware validates room membership, then writes votes to Redis Sorted
Sets via atomic \texttt{ZINCRBY} and prevents double-voting with a Redis Set;
(3) updated tallies are broadcast to all room members; (4) when a majority (or
plurality once everyone has voted) is reached, the winner is persisted to
PostgreSQL and the Redis keys are deleted.

% ===================== 4. DATABASE DESIGN =====================
\section{Database Design}
\subsection{Core Schema}
The schema is normalised to 3NF. Junction tables (\texttt{user\_libraries},
\texttt{room\_members}, \texttt{session\_ratings}) resolve the many-to-many
relationships with composite primary keys.

\begin{lstlisting}[language=SQL]
profiles(id PK -> auth.users, username UNIQUE, avatar_url,
         preferences JSONB, created_at)
games(id PK, title, genre TEXT[], min_players, max_players,
      avg_playtime_minutes, cover_url)
user_libraries(user_id FK, game_id FK, PK(user_id, game_id))
rooms(id PK, name, host_id FK, status, created_at)
room_members(room_id FK, user_id FK, joined_at, PK(room_id, user_id))
voting_sessions(id PK, room_id FK, started_at, ended_at, winner_game_id FK)
game_sessions(id PK, room_id FK, game_id FK, started_at, ended_at,
              duration_minutes, notes)
session_ratings(session_id FK, user_id FK, rating CHECK(1..5),
                PK(session_id, user_id))
\end{lstlisting}

Bonus features add \texttt{room\_messages} (real-time chat), \texttt{game\_wishlists}
(wishlist set queries), and an audit-log table populated by triggers.

\subsection{JSONB Preference Document}
The \texttt{profiles.preferences} column holds a flexible document, indexed with a
GIN index for fast containment queries:
\begin{lstlisting}[language=SQL]
{
  "favorite_genres": ["FPS", "RPG"],
  "avg_session_minutes": 90,
  "genre_weights": { "FPS": 0.8, "RPG": 0.6, "Strategy": 0.3 },
  "total_sessions": 12
}
\end{lstlisting}

% ===================== 5. ADBMS CONCEPTS =====================
\section{Advanced DBMS Concepts Demonstrated}
\begin{longtable}{p{4.2cm} p{9.6cm}}
\toprule
\textbf{Concept} & \textbf{Where / How Used} \\
\midrule
JSONB document model & \texttt{profiles.preferences} --- flexible schema inside PostgreSQL, updated with an exponential moving average after each rating. \\
Set intersection & \texttt{get\_common\_games} finds games every room member owns. \\
Aggregates + HAVING & \texttt{HAVING COUNT(DISTINCT user\_id) = member\_count} guarantees a true intersection. \\
Stored functions & 8 PL/pgSQL functions encapsulate business logic (see \S\ref{sec:impl}). \\
Triggers & Auto-create a profile on signup; audit-log triggers on insert/update; preference update after ratings. \\
GIN index & On \texttt{profiles.preferences} (JSONB) and on full-text \texttt{tsvector} columns. \\
Row-Level Security & Per-user and per-room access control on all tables. \\
In-memory DB (Redis) & Sorted Sets for live vote tallies; Sets to prevent double-voting. \\
Real-time subscriptions & Supabase Realtime drives live room state and chat. \\
Materialized view & Pre-computed room statistics, refreshed concurrently. \\
Window functions & \texttt{DENSE\_RANK() OVER (...)} for a room leaderboard. \\
Full-text search & \texttt{tsvector} + \texttt{websearch\_to\_tsquery} + GIN index for game search. \\
Normalization & All relational tables in 3NF. \\
\bottomrule
\end{longtable}

% ===================== 6. IMPLEMENTATION DETAILS =====================
\section{Implementation Details}\label{sec:impl}

\subsection{Auth and Profile Trigger}
Supabase Auth manages credentials. A PostgreSQL trigger on \texttt{auth.users}
automatically inserts a matching row into \texttt{profiles} on registration, so
the application never has to create profiles manually.

\subsection{Common-Library Intersection}
The signature query of the project --- it returns only games owned by \emph{all}
room members:
\begin{lstlisting}[language=SQL]
SELECT g.*
FROM games g
JOIN user_libraries ul ON ul.game_id = g.id
WHERE ul.user_id IN (SELECT user_id FROM room_members WHERE room_id = $1)
GROUP BY g.id
HAVING COUNT(DISTINCT ul.user_id) =
       (SELECT COUNT(*) FROM room_members WHERE room_id = $1);
\end{lstlisting}

\subsection{Preference-Aware Shortlist}
\texttt{get\_shortlist} ranks the common games by average group preference,
falling back to a neutral 0.5 weight for users with no data yet:
\begin{lstlisting}[language=SQL]
SELECT g.*,
  AVG(COALESCE((p.preferences->'genre_weights'->>g.genre[1])::float, 0.5))
    AS group_score
FROM games g
JOIN user_libraries ul ON ul.game_id = g.id
JOIN profiles p ON p.id = ul.user_id
WHERE ul.user_id IN (SELECT user_id FROM room_members WHERE room_id = $1)
GROUP BY g.id
HAVING COUNT(DISTINCT ul.user_id) =
       (SELECT COUNT(*) FROM room_members WHERE room_id = $1)
ORDER BY group_score DESC
LIMIT 5;
\end{lstlisting}

\subsection{Preference Update (EMA)}
After each session rating, a stored function nudges the user's genre weight using
an exponential moving average, keeping the JSONB profile adaptive:
\begin{lstlisting}[language=SQL]
new_weight = old_weight * 0.8 + (rating / 5.0) * 0.2
\end{lstlisting}

\subsection{Live Voting with Redis}
The middleware models each vote as Redis structures keyed by room and voting
session:
\begin{lstlisting}
vote:{room_id}:{voting_session_id}    -> Sorted Set (member=game_id, score=votes)
voters:{room_id}:{voting_session_id}  -> Set (member=user_id, blocks double-votes)
\end{lstlisting}
On each vote the server checks \texttt{SISMEMBER}, then \texttt{ZINCRBY}, then
\texttt{SADD}, and broadcasts tallies. A game wins on majority
(\texttt{ceil(members/2)}); if all members have voted with no majority, the
plurality leader wins. The winner is persisted to \texttt{voting\_sessions} and a
\texttt{game\_sessions} row is created.

\subsection{Bonus Features}
\begin{itemize}[leftmargin=1.4em]
  \item \textbf{Materialized view} of room statistics with concurrent refresh.
  \item \textbf{Window-function leaderboard} using \texttt{DENSE\_RANK()}.
  \item \textbf{Full-text search} on games via \texttt{tsvector} + GIN + trigger.
  \item \textbf{Audit log} populated by \texttt{AFTER INSERT/UPDATE} triggers with JSONB metadata.
  \item \textbf{Real-time chat} per room via Supabase Realtime with optimistic updates.
  \item \textbf{Rematch / Play Again} and a personal \textbf{Wishlist} with a set-based room wishlist query.
\end{itemize}

% ===================== 7. FEATURES =====================
\section{Features and Functionalities}
\begin{itemize}[leftmargin=1.4em]
  \item \textbf{Authentication:} register, log in, log out, persistent sessions, auto-created profile.
  \item \textbf{Game Library:} browse a seeded catalogue and toggle owned games.
  \item \textbf{Rooms:} create a room (auto invite code), join by code, view members.
  \item \textbf{Common Games:} live intersection of all members' libraries.
  \item \textbf{Smart Shortlist:} top games ranked by aggregate group preference.
  \item \textbf{Live Voting:} real-time vote bars, majority/plurality resolution, no double-voting.
  \item \textbf{Session Logging:} sessions recorded with duration; per-user 1--5 star ratings.
  \item \textbf{Preference Profiles:} JSONB taste profiles that adapt after each rating.
  \item \textbf{History \& Stats:} per-room session history and personal statistics.
  \item \textbf{Bonus:} room chat, leaderboard, full-text game search, audit log, wishlist, rematch.
\end{itemize}

% ===================== 8. DEPLOYMENT =====================
\section{Deployment}
\begin{itemize}[leftmargin=1.4em]
  \item \textbf{Frontend:} deployed on Vercel --- \url{https://gamenight-adbms.vercel.app}
        (publicly available at all times).
  \item \textbf{Database / Auth / Realtime:} hosted on Supabase (cloud).
  \item \textbf{Live-voting middleware:} a Node.js + Redis WebSocket server exposed
        over a secure tunnel during demonstrations; all other features run 24/7.
\end{itemize}

% ===================== 9. RUNNING LOCALLY =====================
\section{Instructions for Running the Project}
\subsection{Prerequisites}
Node.js v18+, Redis, a Supabase project, and Git.

\subsection{Setup}
\begin{lstlisting}[language=bash]
# 1. Clone
git clone https://github.com/izzatkamall/gamenight-adbms.git
cd gamenight-adbms

# 2. Apply database migrations
#    Run every file in supabase/migrations/ in numeric order
#    (001 ... 021) in the Supabase SQL editor.

# 3. Frontend environment  ->  frontend/.env.local
#    VITE_SUPABASE_URL=...
#    VITE_SUPABASE_ANON_KEY=...
#    VITE_MIDDLEWARE_WS_URL=ws://localhost:4000

# 4. Middleware environment ->  middleware/.env
#    REDIS_URL=redis://localhost:6379
#    SUPABASE_URL=...
#    SUPABASE_SERVICE_ROLE_KEY=...
#    PORT=4000
\end{lstlisting}

\begin{lstlisting}[language=bash]
# 5. Install dependencies
cd frontend   && npm install
cd ../middleware && npm install

# 6. Run (three terminals)
redis-server                 # terminal 1
cd middleware && node index.js   # terminal 2
cd frontend  && npm run dev      # terminal 3
# open http://localhost:5173
\end{lstlisting}

% ===================== 10. SCREENSHOTS =====================
\section{Screenshots}

\subsection{Application}
\screenshot{landing.png}{Landing page --- dark, futuristic gaming UI.}{fig:landing}
\screenshot{auth.png}{User authentication (login / register).}{fig:auth}
\screenshot{dashboard.png}{Dashboard with the user's rooms and quick actions.}{fig:dashboard}
\screenshot{library.png}{Game library --- toggling owned games.}{fig:library}
\screenshot{room-lobby.png}{Room lobby --- common games and the preference-ranked shortlist.}{fig:lobby}
\screenshot{voting.png}{Live voting screen with real-time vote tallies.}{fig:voting}
\screenshot{session-history.png}{Session history / results.}{fig:history}
\screenshot{profile.png}{Profile page --- genre preference breakdown and stats.}{fig:profile}

\subsection{Database and Backend}
\screenshot{supabase-tables.png}{Supabase Table Editor --- normalised schema with RLS enabled.}{fig:tables}
\screenshot{sql-function.png}{Stored function returning a room's common games / shortlist.}{fig:func}
\screenshot{jsonb-preferences.png}{A profile's \texttt{preferences} JSONB document with genre weights.}{fig:jsonb}
\screenshot{redis-voting.png}{Redis Sorted Set holding live vote tallies during a session.}{fig:redis}

% ===================== 11. CONCLUSION =====================
\section{Additional Information for Evaluation}
GameNight demonstrates how a single application can responsibly combine a
relational core (normalised schema, stored functions, triggers, RLS) with a
document model (JSONB), an in-memory store (Redis) for a genuinely
latency-sensitive feature, and real-time data flow. Each technology is used where
it is the right tool rather than for its own sake, and the live-voting path in
particular justifies Redis through its atomic, high-frequency write pattern.

The complete source code, all 21 SQL migrations, and commit history are available
in the public repository linked on the title page.

\end{document}
